\begin{theorem}
Let $1<p<\infty$ with $p\neq2$, let $k\ge1$ be an integer and $0<\alpha\le1$
with $k+\alpha$ equal to the right-hand side of (9), as in
\noderef{KAlphaExp}, and let $\Omega\subseteq\mathbb{C}$ be open, non-empty and connected.
Then there is a function $v\colon\mathbb{C}\to\mathbb{R}$ such that
\begin{enumerate}
\item[(a)] $v$ is $p$-harmonic on $\Omega$ in the sense of
  \noderef{PHarmonicOn};
\item[(b)] $v$ is continuous on $\Omega$;
\item[(c)] $v\notin C^{k+\alpha}_{\mathrm{loc}}(\Omega)$ in the sense of
  \noderef{MemHolderSmallOn};
\item[(d)] $v\notin W^{k+2,q}_{\mathrm{loc}}(\Omega)$ in the sense of
  \noderef{MemWLocOn}, where $q=\frac{2}{2-\alpha}$.
\end{enumerate}

This is the optimality clause of Theorem 1 of the paper: ``for each
$1<p<\infty$, $p\neq2$, there is a $p$-harmonic function
$v\in W^{1,p}_{\mathrm{loc}}(\Omega)$ which is not in the class
$C^{k+\alpha}_{\mathrm{loc}}(\Omega)\cup W^{k+2,\frac{2}{2-\alpha}}
_{\mathrm{loc}}(\Omega)$''; the membership $v\in W^{1,p}_{\mathrm{loc}}(\Omega)$
is part of \noderef{PHarmonicOn}.  Together with the regularity half it covers
the target: neither half covers it alone.

Two features of the statement deserve comment.  The conclusion concerns $v$,
whose two indices exceed by one those of (56), which concerns the complex
gradient $f=\partial v/\partial z$; the exponent $q$ does not move.  And the
continuity (b) is asserted deliberately.  Being $p$-harmonic is a property of
$v$ seen only through integrals, so it is unchanged when $v$ is altered on a
set of measure zero, whereas membership in a Hölder class is not; without (b)
the non-membership (c) could be arranged by altering any function on a null
set, and the statement would be vacuous.  It is not a strengthening the paper
does not justify: the paper's own witness, the potential of the continuous
function $f$ of (55), is continuous.  Connectedness of $\Omega$ carries its
non-emptiness, without which (c) and (d) are false, every function lying in
every class on the empty set; and it is the paper's hypothesis on the domain.
\end{theorem}
\begin{proof}
One remark on well-posedness before the argument proper.  The hypothesis above
presents the pair $(k,\alpha)$ as given, whereas Theorem 1 of the paper says at
paper.tex:87--90 that ``the integer $k\ge1$ and the exponent $\alpha\in(0,1]$
are determined by the equation (9)''.  By \noderef{KAlphaWellDefined} there is,
for the given $p$, exactly one pair consisting of an integer $k\ge1$ and a real
number $\alpha$ with $0<\alpha\le1$ such that $k+\alpha$ equals the right-hand
side of (9); so the hypothesis on $(k,\alpha)$ is satisfiable for every
$1<p<\infty$, and in particular for every $p\neq2$, and determines the pair from
$p$, exactly as the paper asserts, and the statement proved below is therefore
not vacuous for any such $p$.  This is a well-posedness remark and not a step in
the argument: nothing in the rest of this proof invokes that lemma, and in
particular the bundle fed to \noderef{Section7Heart} in Step 3 below uses only
the hypothesis on $(k,\alpha)$ as stated, for whichever pair satisfies it.

\emph{Step 1: the extremal data.}  By \noderef{ExtremalInverseGradient} there
are $f\colon\mathbb{C}\to\mathbb{C}$ and $v_0\colon\mathbb{C}\to\mathbb{R}$ such
that: $f$ is continuous on $\mathbb{C}$ and satisfies $f(H(\xi))=\xi$ and
$H(f(z))=z$ for all $\xi,z$, where $H=H_p$ is the extremal function of
\noderef{ExtremalH}; $f$ is $C^\infty$ on $\mathbb{C}-\{0\}$; $f$ satisfies the
two homogeneity identities of degree $d=d(p)$ and $d-|\nu|$ of
\noderef{DExp} and \noderef{PartialWord}; $v_0$ is continuous on $\mathbb{C}$,
is $C^\infty$ on $\mathbb{C}-\{0\}$ and is $p$-harmonic on $\mathbb{C}$ in the
sense of \noderef{PHarmonicOn}; and
\begin{equation}\tag{1}
  \frac{\partial v_0}{\partial z}(z)=f(z)\qquad(z\neq0),
\end{equation}
with $\partial/\partial z$ as in \noderef{DZ}.

\emph{Step 2: translation of the singular point into $\Omega$.}  Since $\Omega$
is connected it is non-empty; fix $z_0\in\Omega$ and put
\[
  U=\{z-z_0:z\in\Omega\},\qquad v(z)=v_0(z-z_0)\quad(z\in\mathbb{C}).
\]
Then $U$ is open, being the image of the open set $\Omega$ under the
homeomorphism $z\mapsto z-z_0$, and $0=z_0-z_0\in U$.  Write
$\tau(z)=z-z_0$, so that $v=v_0\circ\tau$, $U=\tau(\Omega)$ and
$\Omega=\tau^{-1}(U)$.

\emph{Step 3: (56) on $U$.}  The hypotheses of \noderef{Section7Heart} hold with
this $U$ and this $f$: $1<p<\infty$ and $p\neq2$; $k\ge1$, $0<\alpha\le1$ and
$k+\alpha$ is the right-hand side of (9); $U$ is open with $0\in U$; and the
four properties (H1)--(H4) of $f$ --- smoothness off the origin, the two
inverse relations, and the two homogeneity identities --- are precisely those
listed in Step 1.
Therefore
\begin{equation}\tag{2}
  f\notin C^{k-1+\alpha}_{\mathrm{loc}}(U)
  \quad\text{and}\quad
  f\notin W^{k+1,q}_{\mathrm{loc}}(U).
\end{equation}

\emph{Step 4: two remarks on norms and on translation.}

\noindent(4a) Let $\Omega'$ be open, let $g$ be of class $C^m$ on $\Omega'$ and
let $\mathrm{d}^mg$ be its $m$-th iterated Fréchet derivative, which on the
open set $\Omega'$ agrees with the derivative within $\Omega'$.  As in
\noderef{MemHolderSmallOn} the little-o condition is expressed through
$\mathrm{d}^mg$, whereas the partial derivatives $D^\nu g$ of
\noderef{PartialWord} with $|\nu|=m$ are the values
$\mathrm{d}^mg(x)[e_{j_1},\dots,e_{j_m}]$ on the basis vectors $e_0=1$,
$e_1=i$.  Expanding each argument of an $m$-linear map in that basis and
using $|u^{(j)}|\le\|u\|$ gives, for $x,y\in\Omega'$,
\begin{equation}\tag{3}
  \bigl|D^\nu g(x)-D^\nu g(y)\bigr|
    \le\bigl\|\mathrm{d}^mg(x)-\mathrm{d}^mg(y)\bigr\|
    \le\sum_{|\nu'|=m}\bigl|D^{\nu'}g(x)-D^{\nu'}g(y)\bigr| ,
\end{equation}
the last sum having $2^m$ terms.

\noindent(4b) Translation covariance.  Lebesgue measure on $\mathbb{C}$ is
translation invariant, $\tau$ is a $C^\infty$ diffeomorphism with
$\mathrm{d}\tau=\mathrm{id}$, $\tau$ carries compact subsets of $\Omega$ to
compact subsets of $U$ and preserves distances, and
$\phi\mapsto\phi\circ\tau^{-1}$ is a bijection between the test functions for
$\Omega$ of \noderef{IsTestFn} and those for $U$.  Consequently each of the
notions used below is covariant under $\tau$: for $u\colon\mathbb{C}\to\mathbb{C}$,
\begin{itemize}
\item $u\in L^{s}_{\mathrm{loc}}(\Omega)$ if and only if
  $u\circ\tau^{-1}\in L^{s}_{\mathrm{loc}}(U)$, by
  \noderef{MemLpLocOn} and translation invariance of the measure;
\item $w$ is the distributional $j$-th partial of $u$ on $\Omega$ in the sense
  of \noderef{IsWeakPartialOn} if and only if $w\circ\tau^{-1}$ is that of
  $u\circ\tau^{-1}$ on $U$, by the change of variables $\zeta=\tau(z)$ in both
  integrals, which also converts $\phi$ into $\phi\circ\tau^{-1}$;
\item $u\in C^{m+\alpha}_{\mathrm{loc}}(\Omega)$ if and only if
  $u\circ\tau^{-1}\in C^{m+\alpha}_{\mathrm{loc}}(U)$, by
  \noderef{MemHolderSmallOn}: the $C^m$ clause is preserved by composition with
  a diffeomorphism whose differential is the identity, which also gives
  $\mathrm{d}^m(u\circ\tau^{-1})(\zeta)=\mathrm{d}^mu(\tau^{-1}(\zeta))$, and
  the little-o clause then transfers because $\tau^{-1}$ carries compact
  subsets of $U$ to compact subsets of $\Omega$ and preserves distances;
\item $u\in W^{m,s}_{\mathrm{loc}}(\Omega)$ if and only if
  $u\circ\tau^{-1}\in W^{m,s}_{\mathrm{loc}}(U)$, by \noderef{MemWLocOn}: the
  family $(D_w)$ for $u$ on $\Omega$ transforms into the family
  $(D_w\circ\tau^{-1})$ for $u\circ\tau^{-1}$ on $U$ by the two previous
  items;
\item for $h\colon\mathbb{C}\to\mathbb{C}$, the integral identity of clause (c)
  of \noderef{PHarmonicOn} holds with witness $h$ on $\Omega$ if and only if it
  holds with witness $h\circ\tau^{-1}$ on $U$.  Written out, with
  $\nabla\phi=\phi_x+i\phi_y$ and the integrand that \noderef{PHarmonicOn}
  records for clause (c), the two assertions are
  \[
    \int_{\mathbb{C}}|h|^{p-2}
      \operatorname{Re}\bigl(h\,\overline{\nabla\phi}\bigr)\,d\sigma=0
    \quad\text{for every real test function }\phi\text{ for }\Omega
  \]
  and the same with $h$ replaced by $h\circ\tau^{-1}$ and $\phi$ ranging over
  the real test functions for $U$.  Take the direction $U\to\Omega$, which is
  the one Step 5 uses, and assume the second.  Let $\phi$ be a real test
  function for $\Omega$ and put $\psi=\phi\circ\tau^{-1}$, a real test function
  for $U$ by the bijection recorded above.  Substituting $\zeta=\tau(z)$,
  \[
    \int_{\mathbb{C}}|h(z)|^{p-2}
      \operatorname{Re}\bigl(h(z)\,\overline{\nabla\phi(z)}\bigr)\,d\sigma(z)
    =\int_{\mathbb{C}}\bigl|(h\circ\tau^{-1})(\zeta)\bigr|^{p-2}
      \operatorname{Re}\bigl((h\circ\tau^{-1})(\zeta)\,
        \overline{\nabla\psi(\zeta)}\bigr)\,d\sigma(\zeta),
  \]
  by three facts, each already recorded in the preamble to this list: the
  measure is unchanged because Lebesgue measure on $\mathbb{C}$ is translation
  invariant; the witness is carried along because
  $h(z)=(h\circ\tau^{-1})(\tau z)$ by definition; and the gradient is carried
  along because $\mathrm{d}\tau=\mathrm{id}$ gives
  $\nabla\psi(\zeta)=\nabla\phi(\tau^{-1}\zeta)=\nabla\phi(z)$.  The right-hand
  integral vanishes by the assumption on $U$, so the left-hand one does too;
  and since $\phi\mapsto\phi\circ\tau^{-1}$ is a \emph{bijection} onto the real
  test functions for $U$, the quantifier transfers as well.  The converse
  direction is the same computation run through $\tau^{-1}$, which is again a
  translation.  The convention of \noderef{PHarmonicOn} by which $|h|^{p-2}$ is
  read as $0$ where $h$ vanishes causes no difficulty here: the integrand is a
  fixed function of the two values $h(z)$ and $\nabla\phi(z)$, so the
  substitution transports it along with them.
\end{itemize}
Finally $v\circ\tau^{-1}=v_0$ and $\tau(\Omega)=U$.

\emph{Step 5: (a) and (b).}  By Step 1, $v_0$ is $p$-harmonic on $\mathbb{C}$.
Restricting to $\Omega$ is legitimate: by \noderef{PHarmonicOn} with the
witness $g_0$, the two local $L^p$ memberships and the two
integration-by-parts identities and the integral identity are required against
test functions for the ambient set, and every test function for a subset of
$\mathbb{C}$ is one for $\mathbb{C}$, while compact subsets of a subset are
compact subsets of $\mathbb{C}$; so $v_0$ is $p$-harmonic on $U$ with the same
witness.  Now transfer that back to $\Omega$.  By \noderef{PHarmonicOn} what
has to be checked for $v=v_0\circ\tau$ on $\Omega$ with witness
$g=g_0\circ\tau$ is exactly its three clauses, and each is an instance of (4b)
read in the direction $U\to\Omega$ --- that is, right to left in the stated
equivalence, with the function on the $\Omega$ side taken to be $v_0\circ\tau$
or $g_0\circ\tau$, so that composing it with $\tau^{-1}$ returns the function
already known on $U$.  Clause (a), that $v$ and $g$ lie in
$L^p_{\mathrm{loc}}(\Omega)$, is the first item applied once to each of the two
functions.  Clause (b), that $\operatorname{Re}g$ and $\operatorname{Im}g$ are
the distributional $x$- and $y$-partials of $v$ on $\Omega$, is the second item
applied once for each of the two letters, using
$\operatorname{Re}(g_0\circ\tau)=(\operatorname{Re}g_0)\circ\tau$ and likewise
for the imaginary part.  Clause (c), the integral identity, is the fifth item
with $h=g_0\circ\tau$.  All three therefore hold, so $v$ is $p$-harmonic on
$\Omega$, which is (a).  And $v$ is continuous on $\mathbb{C}$, being the
composition of the continuous $v_0$ with $\tau$; in particular it is continuous
on $\Omega$, which is (b).

\emph{Step 6: (c).}  Suppose $v\in C^{k+\alpha}_{\mathrm{loc}}(\Omega)$.  By the
third item of (4b), $v_0=v\circ\tau^{-1}\in C^{k+\alpha}_{\mathrm{loc}}(U)$.
We show that then $f\in C^{k-1+\alpha}_{\mathrm{loc}}(U)$, contradicting the
first half of (2).

By \noderef{MemHolderSmallOn} there are two cases, and in both $v_0$ is at
least of class $C^{k}$ on $U$: if $\alpha=1$ the clause is
$v_0\in C^{k+1}(U)$, and if $\alpha<1$ it is $v_0\in C^{k}(U)$ together with
the little-o condition on $\mathrm{d}^{k}v_0$.  Put
\[
  w_0=\frac{\partial v_0}{\partial z}
  =\tfrac12\bigl(\partial_x v_0-i\,\partial_y v_0\bigr),
\]
which by \noderef{DZ} is the stated combination of first partial derivatives,
and which is therefore of class $C^{k-1}$ on $U$, with
\begin{equation}\tag{4}
  D^\nu w_0=\tfrac12\bigl(D^{\nu\frown(0)}v_0-i\,D^{\nu\frown(1)}v_0\bigr)
  \qquad(|\nu|\le k-1)
\end{equation}
at every point of $U$, where $\nu\frown(j)$ is $\nu$ with the letter $j$
appended on the right, a word of length $|\nu|+1$.  Two things make (4)
legitimate, and both need the set named.  First,
$D^{\nu\frown(j)}=D^{\nu}\circ\partial_j$: by \noderef{PartialWord} the
operator attached to a word is the composite
$\partial_{j_1}\partial_{j_2}\cdots\partial_{j_m}$ in which the last letter is
applied first, so appending $j$ on the right applies $\partial_j$ first.
Second, the step is linearity of $D^{\nu}$, which is \emph{not} available
unconditionally: by the convention of \noderef{PartialWord} a derivative is
taken to be $0$ wherever the function fails to be real-differentiable, so these
operators are additive only where both summands are differentiable.  Here the
instance is located.  The set is the open set $U$, on which $v_0$ is of class
$C^{k}$ as just established; for $|\nu|\le k-1$ each of the at most $k$
applications of $\partial_0$ or $\partial_1$ occurring on either side of (4)
acts on a function of class $C^{k}$ on $U$, so at every point of $U$ all the
derivatives involved are genuine and the convention never intervenes.  Since
$U$ is open, that is all (4) asserts.

First, $w_0=f$ on $U$.  Both are continuous on $U$: $w_0$ because
$v_0\in C^{k}(U)$ with $k\ge1$, and $f$ by Step 1.  They agree on $U-\{0\}$ by
(1).  The singleton $\{0\}$ has empty interior, so $U-\{0\}$ is dense in $U$,
and two continuous functions on $U$ agreeing on a dense subset agree on $U$.

Now the index shift.  If $\alpha=1$, then $v_0\in C^{k+1}(U)$, so
$w_0\in C^{k}(U)$ by (4), and by \noderef{MemHolderSmallOn} at the index $k-1$
and $\alpha=1$ --- whose clause is $C^{(k-1)+1}=C^{k}$, using $k\ge1$ --- we get
$w_0\in C^{k-1+\alpha}_{\mathrm{loc}}(U)$.  If $\alpha<1$, then
$w_0\in C^{k-1}(U)$, and we must verify the little-o clause for $w_0$ at the
index $k-1$.  Let $K\subseteq U$ be compact and $\delta>0$.  Apply the little-o
clause for $v_0$ at the index $k$ with the same $K$ and with
$\delta/2^{k-1}>0$, obtaining $\rho>0$ such that
$\|\mathrm{d}^{k}v_0(x)-\mathrm{d}^{k}v_0(y)\|\le2^{-(k-1)}\delta|x-y|^{\alpha}$
for $x,y\in K$ with $|x-y|<\rho$.  For such $x,y$ and for $|\nu|=k-1$, (4) and
the first inequality of (3) applied to $v_0$ at the order $k$ give
\[
  \bigl|D^\nu w_0(x)-D^\nu w_0(y)\bigr|
  \le\tfrac12\sum_{j\in\{0,1\}}
    \bigl|D^{\nu\frown(j)}v_0(x)-D^{\nu\frown(j)}v_0(y)\bigr|
  \le\bigl\|\mathrm{d}^{k}v_0(x)-\mathrm{d}^{k}v_0(y)\bigr\|
  \le\frac{\delta|x-y|^{\alpha}}{2^{k-1}} ,
\]
and then the second inequality of (3) applied to $w_0$ at the order $k-1$,
whose sum has $2^{k-1}$ terms, gives
\[
  \bigl\|\mathrm{d}^{k-1}w_0(x)-\mathrm{d}^{k-1}w_0(y)\bigr\|
  \le2^{k-1}\cdot\frac{\delta|x-y|^{\alpha}}{2^{k-1}}
  =\delta|x-y|^{\alpha}.
\]
This is the little-o clause for $w_0$, so again
$w_0\in C^{k-1+\alpha}_{\mathrm{loc}}(U)$.

Since $w_0=f$ on $U$, and since both clauses of \noderef{MemHolderSmallOn}
involve the function only through its restriction to $U$ --- the smoothness
being smoothness on $U$ and the derivatives being taken within $U$ --- we
conclude $f\in C^{k-1+\alpha}_{\mathrm{loc}}(U)$, contradicting (2).
Hence (c) holds.

\emph{Step 7: (d).}  Suppose $v\in W^{k+2,q}_{\mathrm{loc}}(\Omega)$.  By the
fourth item of (4b), $v_0\in W^{k+2,q}_{\mathrm{loc}}(U)$.  We show that then
$f\in W^{k+1,q}_{\mathrm{loc}}(U)$, contradicting the second half of (2).

Let $(D_w)_w$ be a family as in \noderef{MemWLocOn} for $v_0$ on $U$ at the
order $k+2$: $D_{()}=v_0$, each $D_w$ with $|w|\le k+2$ lies in
$L^{q}_{\mathrm{loc}}(U)$, and $D_{w\frown(j)}$ is the distributional $j$-th
partial of $D_w$ on $U$ whenever $|w|<k+2$.  Note that
\noderef{IsWeakPartialOn} is an identity between integrals which is linear in
the pair of functions involved: if $w_1$ is the distributional $j$-th partial of
$u_1$ and $w_2$ that of $u_2$, then $c_1w_1+c_2w_2$ is that of
$c_1u_1+c_2u_2$; and it is unchanged if either function is replaced by one
agreeing with it almost everywhere.

Define
\[
  E_{()}=f,\qquad
  E_w=\tfrac12\bigl(D_{(0)\frown w}-i\,D_{(1)\frown w}\bigr)
  \quad\text{for }1\le|w|\le k+1,
\]
where $(j)\frown w$ is the word $w$ with the letter $j$ prepended, of length
$|w|+1\le k+2$.

\noindent(7a) \emph{$E_{()}$ agrees almost everywhere with
$\tfrac12(D_{(0)}-iD_{(1)})$.}  On the open set $U-\{0\}$ the function $v_0$ is
$C^\infty$ by Step 1.  Restricting the integration-by-parts identities to test
functions for $U-\{0\}$, each $D_{(j)}$ is the distributional $j$-th partial of
$v_0$ on $U-\{0\}$.  Apply \noderef{WeakPartialEqFDerivAE} on $U-\{0\}$ to
$v_0$ and the pair $j\mapsto D_{(j)}$: its hypotheses hold, since $v_0$ is
continuous on $U$ and hence in $L^1_{\mathrm{loc}}(U-\{0\})$; each $D_{(j)}$
lies in $L^{q}_{\mathrm{loc}}(U)$ with $q\ge1$, hence in
$L^1_{\mathrm{loc}}(U-\{0\})$ by Hölder's inequality on compact sets, which
have finite measure; the two identities were just noted; and $v_0$ is
differentiable at every point of $U-\{0\}$.  The conclusion is that
$D_{(j)}=\partial_jv_0$ almost everywhere on $U-\{0\}$ for each $j$, so by
\noderef{DZ} and (1),
\[
  \tfrac12\bigl(D_{(0)}-iD_{(1)}\bigr)
  =\tfrac12\bigl(\partial_xv_0-i\,\partial_yv_0\bigr)
  =\frac{\partial v_0}{\partial z}=f=E_{()}
\]
almost everywhere on $U-\{0\}$, hence almost everywhere on $U$, the origin
being a null set.

\noindent(7b) \emph{$(E_w)$ is a family as required.}  We have $E_{()}=f$.  Each
$E_w$ with $1\le|w|\le k+1$ is a linear combination of two members of
$L^{q}_{\mathrm{loc}}(U)$, hence lies in $L^{q}_{\mathrm{loc}}(U)$; and
$E_{()}=f$ is continuous on $U$, hence bounded on compact subsets of $U$ and so
in $L^{q}_{\mathrm{loc}}(U)$ as well.  Let $|w|\le k$ and $j\in\{0,1\}$; we
check that $E_{w\frown(j)}$ is the distributional $j$-th partial of $E_w$ on
$U$.  Since $|(j')\frown w|=|w|+1\le k+1<k+2$ for each letter $j'$, the family
$(D_w)$ gives that $D_{(j')\frown w\frown(j)}$ is the distributional $j$-th
partial of $D_{(j')\frown w}$ on $U$; by linearity, 
\[
  E_{w\frown(j)}
  =\tfrac12\bigl(D_{(0)\frown w\frown(j)}-i\,D_{(1)\frown w\frown(j)}\bigr)
\]
is the distributional $j$-th partial of
$\tfrac12(D_{(0)\frown w}-iD_{(1)\frown w})$ on $U$.  If $|w|\ge1$ the latter is
$E_w$ and we are done; if $w=()$ the latter agrees with $E_{()}=f$ almost
everywhere by (7a), and the integral identity is insensitive to that, so
$E_{(j)}$ is the distributional $j$-th partial of $E_{()}$ as well.

Hence $f\in W^{k+1,q}_{\mathrm{loc}}(U)$ in the sense of \noderef{MemWLocOn},
contradicting (2).  This proves (d), and with Steps 5 and 6 the theorem.
\end{proof}
